% ---------- Tagged PDF (optional) ----------
% Uncomment the next line for a tagged PDF: real reading order and list
% semantics for screen readers. Works with pdflatex on TeX Live 2023+.
% The LaTeX tagging project still labels this experimental, so it is off
% by default.
% \DocumentMetadata{lang=en-US, testphase={phase-III}}

\documentclass[a4paper,10pt]{article}

% lmodern + T1: keeps text extraction intact when tagging (\DocumentMetadata) is enabled
\usepackage{lmodern}
\usepackage[T1]{fontenc}

% -------------------- Packages --------------------
\usepackage[a4paper,margin=0.6in]{geometry}
\usepackage{titlesec}
\usepackage{enumitem}
\usepackage[hidelinks]{hyperref}
\usepackage{etoolbox}

% -------------------- Setup  --------------------
\hypersetup{
    pdftitle={Ludwig Alvarado Becerra - CV},
    pdfauthor={Ludwig Alvarado Becerra},
    pdfsubject={Curriculum Vitae},
    pdfkeywords={Ingeniería de Software, Ciencia de Datos, Modelado Computacional, Linux, Python, C\#, CI/CD, Bash, Git},
    pdfcreator={LaTeX with hyperref}
}

% -------------------- Formatting --------------------
% Deliberately no font package — Computer Modern, the "tech standard" look.
\pagestyle{empty}
\setlength{\parindent}{0pt}
\setlength{\parskip}{4pt}
\linespread{1.05}
\setlist[itemize]{itemsep=-2pt,topsep=2pt,leftmargin=0.18in}
\sloppy
\raggedbottom

% Section formatting — small-caps title, full-width rule immediately under it,
% tight spacing (Jake's Resume's signature look).
\titleformat{\section}
  {\large\scshape\raggedright}
  {}{0em}{}[\titlerule]

\titlespacing*{\section}{0pt}{7pt}{3pt}

% -------------------- Semantic macros --------------------
% render.py (app/services/conversion/render.py) only ever emits calls to the 11
% macros below, with data already massaged (dates joined, "Present" resolved,
% link labels derived, etc.) but with ALL visual decisions — bold/italic, hfill
% alignment, icons, separators, list styling, gap sizes — owned here. See the
% macro contract in that module's docstring before changing signatures.
% Every template MUST define all 11, including \closebullets and \entrygap.

\newtoggle{citemnotfirst}
\newtoggle{bulletsopen}
\newtoggle{skilllinenotfirst}

% \closebullets closes the itemize list \rbullet opens. render.py calls it
% right after emitting the \rbullet run for an entry. \AtEndDocument is a
% defensive backstop only, in case a future call site forgets.
\newcommand{\closebullets}{%
  \iftoggle{bulletsopen}{\end{itemize}\togglefalse{bulletsopen}}{}%
}
\AtEndDocument{\closebullets}

\newcommand{\resumeheader}[3]{%
  \begin{center}
  {\Huge \scshape #1}\\[3pt]
  \ifblank{#2}{}{{\normalsize #2}\\[2pt]}%
  \small
  \togglefalse{citemnotfirst}%
  #3
  \end{center}
  \vspace{-8pt}
}

% ATS-safe by design: icon-free — #1 (kind) never selects a glyph. It is only
% checked once, for "email", so that case gets a real mailto: link instead of
% a broken relative-URL href; every other kind renders identically, separated
% by a plain " | " on one centered line.
\newcommand{\citem}[3]{%
  \iftoggle{citemnotfirst}{ $|$ }{\toggletrue{citemnotfirst}}%
  \ifstrequal{#1}{email}{\href{mailto:#2}{#3}}{\ifblank{#2}{#3}{\href{#2}{#3}}}%
}

\newcommand{\sectionstart}[1]{\togglefalse{skilllinenotfirst}\section*{#1}}

% Inverse of Classic: title leads (bold, dates right-flushed on the same
% line); company/location follow on the second line, both italic.
\newcommand{\job}[4]{%
  \textbf{#1}\ifblank{#4}{}{ \hfill #4}\\
  \textit{#2}\ifblank{#3}{}{ \hfill \textit{#3}}%
}

\newcommand{\rbullet}[1]{%
  \iftoggle{bulletsopen}{}{\begin{itemize}\toggletrue{bulletsopen}}%
  \item #1%
}

\newcommand{\skillline}[2]{%
  \iftoggle{skilllinenotfirst}{\\}{\toggletrue{skilllinenotfirst}}%
  \textbf{#1:} #2%
}

\newcommand{\edu}[4]{%
  \textbf{#1}\ifblank{#3}{}{ \hfill #3}\\
  #2\ifblank{#4}{}{ \hfill GPA: #4}%
}

\newcommand{\cert}[1]{\\ \textbf{Certification:} #1}

\newcommand{\project}[4]{%
  \textbf{#1}\ifblank{#2}{}{ -- #2}\ifblank{#3}{}{ \hfill \href{#3}{[#4]}}%
}

\newcommand{\entrygap}{\vspace{4pt}}

% -------------------- Document --------------------

\begin{document}

\resumeheader{Ludwig Alvarado Becerra}{Ingeniero de Sistemas \textbar{} Modelado \& Simulación Computacional}{%
\citem{email}{me@ludwigalvarado.me}{me@ludwigalvarado.me}
\citem{location}{}{Bogotá, Colombia}
\citem{web}{https://ludwigalvarado.me}{www.ludwigalvarado.me}
\citem{github}{https://github.com/theludway}{github/theludway}
\citem{linkedin}{https://www.linkedin.com/in/theludway}{linkedin/theludway}}

\sectionstart{PERFIL}
Ingeniero de Sistemas y doble titulación en Modelado \& Simulación Computacional, con experiencia en desarrollo de software, automatización y mejora de procesos de ingeniería. Experiencia trabajando con sistemas GNU/Linux, CI/CD y desarrollo de aplicaciones, combinando fundamentos de programación, modelado matemático y análisis computacional para resolver problemas técnicos en entornos académicos e industriales.


\sectionstart{EXPERIENCIA}
\job{Practicante Ingeniero de Pruebas}{ROSEN - RTRC Bogotá}{Bogotá, Colombia}{Ene. 2026 -- Jul. 2026}
\rbullet{Automatización de tareas repetitivas en inspección de oleoductos, con impacto estimado de ahorro de \(\sim\)1.5 millones de USD en 4 meses.}
\rbullet{Mejoras de rendimiento en CI/CD reduciendo en 30\% el tiempo de construcción de aplicaciones.}
\rbullet{Migración de módulos legado a aplicaciones modernas en C\# y validación funcional con usuarios y stakeholders.}
\closebullets
\entrygap

\job{Miembro fundador}{Asociación de Cartografía Colaborativa de Colombia (AC3)}{Bogotá, Colombia}{Nov. 2024 -- Presente}
\rbullet{Coordinación técnica de captura de imágenes satelitales con drones usando Docker y GNU/Linux.}
\rbullet{Automatización de instalación y configuración de sistemas Linux mediante Bash.}
\closebullets

\job{Profesor Asistente}{Universidad Jorge Tadeo Lozano}{Bogotá, Colombia}{Feb. 2025 -- Presente}
\rbullet{Diseño de talleres y evaluaciones, apoyo en algoritmos, estructuras de datos y programación competitiva.}
\rbullet{Fomento de participación estudiantil en semilleros y competencias como ICPC y CCPL.}
\closebullets
\entrygap

\sectionstart{EDUCACIÓN}
\job{Ingeniería de Sistemas}{Universidad Jorge Tadeo Lozano}{Bogotá, Colombia}{Ene. 2022 -- Nov. 2026}
\rbullet{Líder del semillero Segmentation Fault, orientado en la programación competitiva y hackathones.}
\rbullet{Fundador de GNUTADEO, colectivo de software libre enfocado en la divulgación de GNU/Linux.}
\closebullets
\entrygap
\job{Modelado \& Simulación Computacional}{Universidad Jorge Tadeo Lozano}{Bogotá, Colombia}{Ene. 2024 -- Nov. 2026}
\rbullet{Doble titulación enfocada en simulación, análisis computacional y matemática aplicada.}
\closebullets
\entrygap

\sectionstart{RECONOCIMIENTOS}
\job{Primer puesto Hackathon RiskTech}{RiskTech}{Bogotá, Colombia}{Nov. 2025}
\rbullet{Diseño de un flujo de trabajo para la detección de fraude en transacciones bancarias utilizando aprendizaje de máquina.}
\closebullets
\entrygap
\job{Segundo puesto Hackathon Desafío Inteligente}{Ministerio de Minas y Energía -- Universidad Tecnológica de Pereira}{Bogotá, Colombia}{Dic. 2025}
\rbullet{Software de gemelos digitales para gestionar los niveles de embalses en Colombia mediante técnicas de modelado y replicabilidad por software libre.}

\sectionstart{HABILIDADES}
\skillline{Tecnologías}{C\#, Python, C++, Java, SQL, Bash, \LaTeX, JavaScript, TypeScript.}
\skillline{DevOps}{CI/CD, automatización, GitHub Actions, GNU/Linux, Git, Docker, Podman, Kubernetes.}
\skillline{Idiomas}{Español nativo, Inglés (B2), Mandarín (HSK3).}

\end{document}
